%!TeX root=MemoriaTFG.tex

\chapter{Introducción}
\label{cap:intro}

En el presente Trabajo Final de Grado (TFG) se evalúa de forma
sistemática el rendimiento de los modelos fundacionales disponibles
para imagen dermatológica clínica, sobre un conjunto homogéneo de
protocolos experimentales y datasets externos públicos. El trabajo
se desarrolla en el marco de dos colaboraciones preexistentes con
la Dra.~R.~Taberner del Hospital Universitario Son Llàtzer (HUSLL)
y se acompaña de la construcción de un dataset propio en español,
DermapixelAI~1.0, que se entrega como contribución independiente.

En este capítulo se presenta el contexto clínico del problema, el
marco operativo en el que se inscribe el trabajo, la pregunta de
investigación que articula el estudio, los objetivos formales, las
restricciones explícitas del diseño y la estructura del resto del
documento. El estado del arte específico ---modelos fundacionales
para imagen médica y dermatológica, métodos de adaptación y
evaluación zero-shot--- se desarrolla en el
capítulo~\ref{cap:sota}.


\section{Contexto clínico}
\label{sec:contexto}

La dermatología clínica atiende un espectro amplio de entidades con
distribución concentrada en un número reducido de motivos de
consulta. El estudio observacional DIADERM, sobre $10\,999$
diagnósticos en $8\,953$ pacientes en territorio español,
ordena los diez diagnósticos más prevalentes como queratosis
actínica, carcinoma basocelular, nevus melanocítico, queratosis
seborreica, otras neoplasias benignas, psoriasis, acné, verrugas
víricas, otros trastornos de la pigmentación y otras
dermatitis~\cite{diaderm2018}. El trabajo local de Taberner et~al.\
sobre el servicio dermatológico del HUSLL documenta $213$
diagnósticos diferentes en un único año
asistencial~\cite{taberner2010_motivos}. Ambos coinciden en un peso
elevado de dermatosis benignas, inflamatorias e infecciosas
habituales.

La sensibilidad temporal del proceso diagnóstico resulta crítica en
el melanoma cutáneo. La supervivencia específica a cinco años
desciende desde aproximadamente $99\%$ en estadio~I hasta
$25\%$ en estadio~IV según la octava edición de la American
Joint Committee on Cancer~\cite{ajcc8}. El acceso al especialista
se realiza mediante derivación desde atención primaria. La
capacidad asistencial del circuito presencial es limitada frente a
una demanda creciente, lo que justifica el interés por sistemas de
apoyo al diagnóstico basados en imagen.

El sustrato visual del diagnóstico dermatológico se articula en tres
modalidades: fotografía clínica (presentación general y contexto
anatómico), dermatoscopia (patrones pigmentarios y estructurales no
visibles a simple vista) e histopatología (confirmación
tisular). Sobre esta tríada operan los sistemas de apoyo basados en
aprendizaje automático.


\section{Marco operativo previo}
\label{sec:trabajo-previo}

El presente trabajo se desarrolla en el marco de dos colaboraciones
preexistentes con la Dra.~R.~Taberner, dermatóloga del HUSLL. La
primera es la modernización tecnológica del archivo Dermapixel,
blog mantenido por la Dra.~Taberner desde 2010 que reúne más de
$700$ casos clínicos comentados con imagen clínica y
dermatoscópica, historia abreviada, diagnóstico y razonamiento
diferencial; en paralelo a este TFG hemos desarrollado su migración a
una pila tecnológica propia. La segunda es un proyecto institucional del
Servei de Salut de les Illes Balears (IB-Salut) sobre
teledermatología hospitalaria, en el que se desarrolla una
herramienta de apoyo a la primera valoración dermatológica en el
circuito de derivación desde atención primaria. Ambas líneas
orientan la decisión de priorizar modelos con código y pesos
públicos por su viabilidad de integración.


\section{Pregunta de investigación}
\label{sec:pregunta}

El estudio se articula en torno a una pregunta única: \emph{¿cuál
es el rendimiento y la capacidad de generalización de los modelos
fundacionales para imagen dermatológica disponibles, en función del
volumen de datos etiquetados y del tipo de supervisión, y con qué
robustez frente a la heterogeneidad taxonómica entre datasets y la
disparidad por fototipo cutáneo?}

La respuesta empírica exige dos condiciones operativas: un
protocolo de evaluación homogéneo aplicable a familias de modelos
heterogéneas y un esquema explícito de armonización taxonómica
entre datasets. El diseño descrito en el
capítulo~\ref{cap:metodos} satisface ambas.


\section{Objetivos del trabajo}
\label{sec:objetivos}

El trabajo se articula en torno a dos objetivos formales,
verificables sobre los resultados experimentales del
capítulo~\ref{cap:resultados}.

\paragraph{O1. Evaluación experimental sistemática de los modelos
fundacionales dermatológicos sobre conjuntos de datos externos.}
Sobre el hardware del proyecto (NVIDIA DGX Spark, chip Grace Hopper
GB10, arquitectura aarch64, $128$\,GB memoria unificada,
CUDA~13.0, BF16) se implementa el conjunto de modelos fundacionales
con código y pesos públicos a la fecha de inicio. Se evalúa su
comportamiento sobre doce datasets externos
públicos\footnote{HAM10000~\cite{ham10000},
BCN20000~\cite{bcn20000}, PAD-UFES-20~\cite{padufes},
DDI~\cite{ddi}, Dermnet~\cite{dermnet}, Derm7pt
(clínica y dermatoscópica)~\cite{kawahara2019}, HIBA, MSKCC, WSI
patches, Fitzpatrick17k~\cite{fitzpatrick17k},
SkinCon~\cite{skincon} e ISIC2018~\cite{isic2018}.} mediante cuatro
protocolos principales (sondeo lineal, ajuste fino supervisado,
segmentación binaria con LoRA~\cite{lora} sobre SAM2.1~\cite{sam2},
clasificación zero-shot multimodal) y tres protocolos
complementarios (comparación con modelos de lenguaje multimodales,
análisis de equidad por fototipo Fitzpatrick, auditoría de
solapamiento con Derm1M mediante hash MD5).

\paragraph{O2. Construcción del dataset DermapixelAI 1.0.}
Se extrae, depura y estructura el material del archivo Dermapixel
para conformar un dataset que vincula, para cada caso, imagen y
modalidad, metadatos clínicos y texto narrativo en español. El
dataset se publica bajo licencia CC-BY-NC-SA~4.0 con datasheet
formal~\cite{datasheets2018}, particiones documentadas y mecanismo
de cita recomendada. La explotación experimental sistemática sobre
el dataset se inicia hacia el cierre del periodo cubierto por este
documento y se desarrolla principalmente con posterioridad;
la secci\'on~\ref{sec:lineas-abiertas} del capítulo~\ref{cap:conclusiones}
recoge el estado del dataset.

Los bloques derivados emergidos durante el desarrollo del proyecto
(armonización ontológica jerárquica L1/L2/L3, ensemble de seguridad
para detección de melanoma, análisis de equidad por fototipo,
interpretabilidad mediante Sparse Autoencoders~\cite{templeton2024},
recuperación visual densa basada en FAISS~\cite{faiss2019},
integración en el prototipo DermApIxel) no constituyen objetivos
formales y se documentan en los anexos correspondientes con su
estado, prerrequisitos y líneas futuras.


\section{Restricciones del estudio}
\label{sec:alcance}

Tres restricciones explícitas condicionan el diseño experimental y
la interpretación de los resultados. \textbf{Hardware:} todas las
ejecuciones se realizan sobre NVIDIA DGX Spark con chip Grace
Hopper GB10 y arquitectura aarch64, no coincidente con el hardware
x86\_64 dominante en la literatura. \textbf{Acceso a datos:} los
conjuntos internos del grupo de Monash empleados en el
preentrenamiento de PanDerm no son accesibles para terceros, lo que
limita el trabajo a los pesos liberados y a la evaluación sobre
datasets externos públicos. \textbf{Acceso a modelos:} los pesos
de DermFM-Zero no son públicos en la fecha de cierre, lo que impide
su evaluación empírica directa y restringe su tratamiento a la
discusión conceptual.

El trabajo no aborda, por decisión de alcance, los siguientes
elementos: preentrenamiento desde cero de un modelo fundacional
propio; validación clínica prospectiva sobre cohorte hospitalaria
real con protocolo aprobado por Comité Ético de Investigación
Clínica; integración operativa con el sistema PACS del HUSLL;
reentrenamiento multilingüe de la rama lingüística sobre datasets en
español de gran escala. Estos elementos se identifican como líneas
futuras en la secci\'on~\ref{sec:lineas-abiertas} del
capítulo~\ref{cap:conclusiones}.


\section{Estructura del documento}
\label{sec:estructura}

El documento se compone de siete capítulos en el cuerpo principal,
una bibliografía y ocho anexos. El capítulo~\ref{cap:sota} sintetiza
el estado del arte en cuatro bloques: aprendizaje profundo aplicado
a dermatología, modelos fundacionales en imagen médica, modelos
específicos para dermatología y brechas operativas no resueltas por
la literatura. El capítulo~\ref{cap:metodos} describe materiales y
métodos: los doce datasets externos, la ontología jerárquica
L1/L2/L3, el dataset DermapixelAI~1.0 (resumido; pipeline detallado
en el Anexo~\ref{cap:anexoc}), los catorce modelos evaluados, la
infraestructura, los protocolos, las métricas y la trazabilidad. El
capítulo~\ref{cap:resultados} reúne los resultados experimentales
organizados por bloque. El capítulo~\ref{cap:discusion} interpreta
los resultados en clave de hallazgos, comparación con la literatura
e implicaciones operativas. El capítulo~\ref{cap:limitaciones}
declara las limitaciones del estudio en seis dimensiones. El
capítulo~\ref{cap:conclusiones} responde de forma compacta a la
pregunta de investigación y enumera las contribuciones
reproducibles.

Los anexos contienen, respectivamente, el mapa de tareas
experimentales y la estructura del repositorio
(Anexo~\ref{cap:anexoa}), las tablas detalladas de resultados por
dataset (Anexo~\ref{cap:anexob}), el pipeline de construcción y
caracterización del dataset DermapixelAI~1.0
(Anexo~\ref{cap:anexoc}), la declaración sobre el uso de
herramientas de inteligencia artificial durante la elaboración del
documento (Anexo~\ref{cap:anexod}), el documento de entrega formal
del dataset con licencia, datasheet y consideraciones de privacidad
(Anexo~\ref{cap:anexoe}), la documentación de los Sparse
Autoencoders y el diccionario de conceptos
(Anexo~\ref{cap:anexof}), la comparación extendida con modelos
de lenguaje multimodales generalistas (Anexo~\ref{cap:anexog}) y
el prototipo DermApIxel (Anexo~\ref{cap:anexoh}).
